\label{sec:architecture}

\subsection{Public API}

The primary public interface is:

\begin{verbatim}
/// Apportion `house_size` seats among states by population.
pub fn apportion(
    populations: &[(String, u64)],
    house_size: usize,
    method: ApportionMethod,
) -> Vec<(String, usize)>

pub enum ApportionMethod {
    HuntingtonHill,
    Webster,
    Adams,
    Jefferson,
    Hamilton,
}
\end{verbatim}

The function returns a vector of (state name, seat count) pairs,
sorted alphabetically by state name.
Total seat count is guaranteed to equal \texttt{house\_size}.
Each state is guaranteed at least 1 seat (constitutional floor).

Secondary utilities:

\begin{verbatim}
/// Check if Hamilton produces the Alabama paradox at this H.
pub fn check_alabama_paradox(
    populations: &[(String, u64)],
    house_size: usize,
) -> Option<(String, usize, usize)>

/// Verify SHA-256 of output against committed reference hash.
pub fn verify_sha256(
    allocation: &[(String, usize)],
    reference_hash: &str,
) -> bool
\end{verbatim}

\subsection{Priority Queue Implementation}

The divisor methods (HH, Webster, Adams, Jefferson) share a
priority-queue core:

\begin{verbatim}
fn apportion_divisor(
    populations: &[(String, u64)],
    house_size: usize,
    priority: fn(pop: u128, seats: u32) -> Priority,
) -> Vec<(String, usize)> {
    let n = populations.len();
    // Step 1: assign constitutional floor (1 seat per state)
    let mut seats = vec![1usize; n];
    // Step 2: build max-heap with priority for (n+1)th seat
    let mut heap = BinaryHeap::new();
    for (i, (_, pop)) in populations.iter().enumerate() {
        heap.push(PriorityEntry {
            priority: priority(*pop as u128, 1),
            index: i,
        });
    }
    // Step 3: award remaining H - n seats
    for _ in n..house_size {
        let top = heap.pop().unwrap();
        seats[top.index] += 1;
        heap.push(PriorityEntry {
            priority: priority(
                populations[top.index].1 as u128,
                seats[top.index] as u32,
            ),
            index: top.index,
        });
    }
    // Return sorted result
    populations.iter().zip(seats.iter())
        .map(|((name, _), &s)| (name.clone(), s))
        .collect()
}
\end{verbatim}

\subsection{Priority Representation}

The \texttt{Priority} type is a rational number stored as two \texttt{u128}
integers:

\begin{verbatim}
struct Priority {
    numerator: u128,
    denominator: u128,
}

impl PartialOrd for Priority {
    fn partial_cmp(&self, other: &Self) -> Option<Ordering> {
        // Cross-multiply to compare fractions exactly:
        // a/b > c/d iff a*d > b*c (for positive b,d)
        (self.numerator * other.denominator)
            .partial_cmp(&(other.numerator * self.denominator))
    }
}
\end{verbatim}

This representation enables exact comparison without floating-point
approximation, which is essential for correct tie-breaking.

\subsection{Hamilton Implementation}

Hamilton does not use the priority queue.
It is implemented separately:

\begin{verbatim}
fn apportion_hamilton(
    populations: &[(String, u64)],
    house_size: usize,
) -> Vec<(String, usize)> {
    let total: u64 = populations.iter().map(|(_, p)| p).sum();
    // Compute exact quotas as rationals
    let quotas: Vec<(u128, u128)> = populations.iter()
        .map(|(_, &p)| (p as u128 * house_size as u128, total as u128))
        .collect();
    // Assign lower quotas
    let mut seats: Vec<usize> = quotas.iter()
        .map(|&(n, d)| (n / d) as usize)
        .collect();
    // Apply constitutional floor
    for s in seats.iter_mut() { *s = (*s).max(1); }
    // Distribute remainder seats by fractional part
    let remainder = house_size.saturating_sub(seats.iter().sum());
    let mut fracs: Vec<(u128, usize)> = quotas.iter()
        .enumerate()
        .map(|(i, &(n, d))| (n % d, i))
        .collect();
    fracs.sort_by(|a, b| b.0.cmp(&a.0));
    for k in 0..remainder {
        seats[fracs[k].1] += 1;
    }
    populations.iter().zip(seats.iter())
        .map(|((name, _), &s)| (name.clone(), s))
        .collect()
}
\end{verbatim}
